\documentclass[../sablona-main.tex]{subfiles}

\begin{document}
\section{Teoretické vlastnosti modelu a diskuse řešení}

V této části jsou rozebrány základní matematické vlastnosti modelu vycházející z grafického znázornění na obrázku \ref{fig: Graficke reseni} a teorie lineárního programování.

\subsection*{Konvexita a charakteristika množiny přípustných řešení}
Množina přípustných řešení $K$, graficky znázorněná jako zelená oblast, je definována průnikem konečného počtu polorovin určených lineárními nerovnicemi. Z hlediska geometrického se jedná o \emph{konvexní polyedr}. Tato vlastnost je klíčová, neboť v konvexní množině platí, že libovolná úsečka spojující dva body množiny leží celým svým průběhem uvnitř této množiny. Z hlediska optimalizace konvexita zaručuje, že každé lokální optimum je zároveň optimem globálním.

\subsection*{Existence řešení}
Vzhledem k charakteru omezení je množina přípustných řešení $K$:
\begin{itemize}
	\item \emph{Neprázdná} – existuje alespoň jeden bod vyhovující všem omezujícím podmínkám.
	\item \emph{Omezená a uzavřená (kompaktní)} – v nákresně je ohraničena ze všech stran.
\end{itemize}
Podle \emph{Weierstrassovy věty} platí, že spojitá funkce (zde lineární účelová funkce) na kompaktní množině vždy nabývá svého minima i maxima. Existence řešení je tedy matematicky zaručena.

\subsection*{Analýza extrémů}
Účelová funkce, v grafu reprezentovaná červenou přímkou (izokvanta), je lineární. Optimální řešení úlohy lineárního programování se vždy nachází v \emph{extrémním bodě} (vrcholu) konvexního polyedru. 

Při maximalizaci posouváme izokvantu rovnoběžně ve směru růstu funkce, dokud neopustí množinu $K$. Bodem dotyku je v tomto modelu vrchol $[30, 30]$, který představuje optimální řešení úlohy.

\subsection*{Shrnutí vlastností}
\begin{itemize}
	\item Model je lineární a splňuje podmínky konvexity.
	\item Řešení existuje a je jednoznačné (vzhledem k odlišnému sklonu účelové funkce a omezujících přímek).
	\item Množina přípustných řešení tvoří kompaktní konvexní obal svých vrcholů.
\end{itemize}
\end{document}